\documentclass[12pt]{article}
\usepackage[utf8]{inputenc}
\usepackage[T1]{fontenc}
\usepackage{hyperref}
\title{Tai D\'on Language Resources for K-12 Curriculum Development}
\author{Compiled by Claude Sonnet 5 (claude-sonnet-5) for Andr\'as Kornai}
\date{}
\begin{document}
\maketitle

\section{Introduction}
Tai D\'on (ISO 639-3 twh), also called White Tai, is a Kra--Dai language of Vietnam, Laos and China. It is related to Tai Dam (Black Tai, blt) but is not mutually intelligible with it, and many Tai D\'on speakers also know Tai Dam. It is written in the Tai Viet script. Search results gave little that is specific to Tai D\'on and nothing designed for schools. Several results are about Tai Dam, and I mark those as such.

\section{Primary Educational Resources}
No Tai D\'on primers, textbooks or school programmes were found in search results.

\section{Linguistic and Reference Materials}
\begin{itemize}
\item Cocoon (CRDO) archive page for Tai D\'on, with a bibliography of early works including Savina's 1910 dictionary and Donaldson's 1976 White Tai lessons: \url{https://cocoon.huma-num.fr/exist/crdo/lexvo.xql?code=twh}
\item Omniglot, Tai D\'on writing and sample text: \url{https://omniglot.com/writing/taidon.htm}
\item Wikipedia, ISO 639 entry for Tai D\'on: \url{https://wikipedia.com/wiki/ISO_639:twh}
\end{itemize}

\section{Religious and Community Resources}
Joshua Project's Tai Dam in Laos page (this is Tai Dam, not Tai D\'on, so check it before use) mentions audio recordings and a New Testament: \url{https://joshuaproject.net/people_groups/print/15187/LA}

\section{Digital and Audio Resources}
Endangered Alphabets profile of the related Tai Dam script (also Tai Viet): \url{https://www.endangeredalphabets.net/?p=420}. Omniglot's Tai Dam page: \url{https://omniglot.com/writing/taidam.htm}.

\section{Curriculum Development Context}
Few scholars read the script fluently, and published teaching material is scarce. Early dictionaries and lessons exist but are old. Curriculum work would build on Tai Viet script materials and new community-based material, and it should keep Tai D\'on distinct from Tai Dam.

\end{document}